\subsection{Open-Meteo als externe Wetterdatenquelle}
Eine häufig genutzte Quelle für Wetterdaten ist die Open-Meteo-API, welche als frei zugängliche und quelloffene Web-\gls{api} zur Bereitstellung meteorologischer Informationen dient. 
Sie ermöglicht den Zugriff auf Wetterdaten ohne Authentifizierung und eignet sich insbesondere für nicht-kommerzielle Anwendungen~\cite{openmeteo}.

Die \gls{api} basiert auf dem \gls{http}-Protokoll und liefert die angeforderten Daten in einem strukturierten \gls{json}-Format. 
Dadurch wird eine einfache Integration in eingebettete Systeme sowie in verschiedene Programmiersprachen ermöglicht. 
Die Anfrage erfolgt über eine \gls{url} mit spezifischen Parametern, beispielsweise zur Angabe von geografischer Breite und Länge. 
Die bereitgestellten Daten umfassen unter anderem Temperatur, Windgeschwindigkeit und Luftfeuchtigkeit, welche sowohl für aktuelle Zustände als auch in Form zeitlich aufgelöster Verlaufsdaten verfügbar sind~\cite{openmeteo}.

Die zugrunde liegenden Wetterinformationen basieren auf einer Kombination verschiedener Datenquellen, darunter Messungen von Satelliten, Radarsystemen sowie weiteren Sensornetzwerken~\cite{openmeteo}. 
Zur Gewährleistung einer hohen Aktualität werden die verwendeten Wettermodelle in kurzen Zeitintervallen aktualisiert.
Durch die standardisierte Bereitstellung über \gls{http} und \gls{json} sowie die kontinuierliche Aktualisierung der Daten stellt die Open-Meteo-\gls{api} eine geeignete Grundlage für die Integration externer Wetterinformationen für den vorliegenden Anwendungsfall dar~\cite{openmeteo}.

Die Kommunikation mit der API erfolgt über \gls{https}-basierte \gls{rest}-Schnittstellen. 
\gls{https} stellt die abgesicherte Variante des HTTP-Protokolls dar und basiert auf der Übertragung von HTTP über \gls{ssl} bzw. \gls{tls}. 
Dabei werden die Kommunikationspartner authentifiziert und die anschließende Datenübertragung vertraulich umgesetzt. 
Dies erhöht die Sicherheit der Kommunikation, kann jedoch zusätzliche Aufwände durch Handshake, Zertifikatsprüfung und kryptografische Operationen verursachen~\cite{naylor2014https}.

Eine \gls{rest}-Schnittstelle ist eine spezialisierte Form einer Programmierschnittstelle, die auf dem Architekturprinzip \gls{rest} basiert.
Sie ermöglicht die Kommunikation zwischen Client und Server über standardisierte HTTP-Methoden sowie eindeutig adressierbare Ressourcen.
Dabei erfolgt der Datenaustausch typischerweise zustandslos und unter Nutzung einheitlicher Schnittstellen, häufig im JSON-Format~\cite{RESTAPI}.